\documentclass[10pt,a4paper]{article}
\usepackage[margin=18mm,top=17mm,bottom=19mm]{geometry}
\usepackage{amsmath,amssymb,graphicx,array,fancyhdr,lastpage,textcomp}
\usepackage{fontspec}
\setmainfont{texgyreheros-regular.otf}[BoldFont=texgyreheros-bold.otf,ItalicFont=texgyreheros-italic.otf,BoldItalicFont=texgyreheros-bolditalic.otf]
\setlength{\parindent}{0pt}
\setlength{\parskip}{3pt}
\pagestyle{fancy}\fancyhf{}
\renewcommand{\headrulewidth}{0pt}\renewcommand{\footrulewidth}{0.3pt}
\fancyfoot[L]{\footnotesize by paper.shrit.in\quad |\quad normal}
\fancyfoot[R]{\footnotesize Page \thepage\ of \pageref{LastPage}}
\newcommand{\question}[3]{\par\vspace{6pt}\noindent\begin{minipage}[t]{0.075\linewidth}\textbf{#1)}\end{minipage}\begin{minipage}[t]{0.855\linewidth}#3\end{minipage}\hfill\begin{minipage}[t]{0.055\linewidth}\raggedleft(#2)\end{minipage}\par}
\begin{document}
{\LARGE\bfseries Question Paper}\hfill {\small Registration No.: \rule{33mm}{0.3pt}}\par\vspace{8pt}
\begin{center}{\large\bfseries MANIPAL FORMAT | PRACTICE PAPER}\\[6pt]{\bfseries THIRD SEMESTER B.TECH. | MIDSEM PRACTICE}\\[4pt]{\bfseries DATA STRUCTURES [CSS 2101]}\\[4pt]{\small COMPUTER SCIENCE AND ENGINEERING}\\[3pt]{\small SET A: ENSEMBLE OF PAST-PAPER QUESTIONS}\end{center}
\textbf{Marks: 30}\hfill\textbf{Suggested duration: 90 minutes}\par\vspace{3pt}\hrule\vspace{5pt}
{\small\textbf{Answer all questions.} Questions 1--5 are MCQs worth 1 mark each. Select one correct option per MCQ. The remaining questions are theory questions worth 25 marks in total; show working and draw labelled diagrams where appropriate. Modules 1-2; module 3 restricted to stacks. Independent practice paper; not an official university examination.}\par
\medskip\textbf{Section A: Multiple-choice questions (5 marks)}\par
\question{1}{1}{Which data structure directly supports undoing the most recent action first?\par (A) Queue\par (B) Stack\par (C) Sorted array\par (D) Adjacency matrix}
\question{2}{1}{Edits A, B and C are pushed in that order onto an undo stack. Which edit does the first undo remove?\par (A) A\par (B) B\par (C) All three\par (D) C}
\question{3}{1}{After edits A, B and C, two undo operations transfer actions to a redo stack. What is the redo stack from bottom to top?\par (A) C, B\par (B) B, C\par (C) A, B\par (D) A, C}
\question{4}{1}{The redo stack contains C, B from bottom to top. Which edit is reapplied by one redo operation?\par (A) A\par (B) C\par (C) B\par (D) Both B and C}
\question{5}{1}{In a conventional two-stack editor, what happens to the redo stack when a new edit is made after an undo?\par (A) It is reversed\par (B) It is cleared\par (C) It is unchanged\par (D) It becomes the undo stack}
\newpage\textbf{Section B: Theory questions (25 marks)}\par
\question{6}{2}{Write a C program to read and display employee details including the date of joining using nested structure. Consider date of joining as nested structure.}
\question{7}{2}{Assume that you are storing marks of students in 3 subjects. Write a C program that: i. Takes the number of students as input. ii. Dynamically allocates a 2D array using pointers. iii. Accepts marks for each subject.Calculates and prints the average marks of each student using only pointer arithmetic.}
\question{8}{2}{Write a C program to dynamically allocate memory for 2D string arrays (array of strings) and perform concatenation using pointers.}
\question{9}{5}{A music streaming service keeps track of songs in a playlist using a singly linked list. Each node stores the following details: SongID (integer), Title (string), Duration (in seconds). Write a C program to perform the following operations: Insert songs into the playlist in the order they are added. Remove all songs whose duration is less than a user-specified value (Use head pointer for performing this operation). Display the updated playlist. Include a main() function to demonstrate all functionalities.}
\question{10}{5}{A mobile app named "MIT-BLR Connect" maintains contact details of alumni in three different modules: i) Add Recent Contacts Module: Displays recently interacted alumni in order of contact - new ones added at the beginning. ii) Traverse Module: Allows viewing alumni details in forward direction. iii) Develop a module to print the data in reverse order. Develop C functions considering singly linked list as the suitable data structure}
\question{11}{5}{Write a function to convert an infix expression to postfix. Include the following functions as per the prototype. i) int InStackPrecedence(char op) ; to return the precedence of operator within the stack ii) int IncomingPrecedence(char op); to return the precedence of the incoming operator. iii) void Infix\_Postfix(char infix[], char postfix[]); to convert given infix expression to postfix. Main function not required. Note: The infix expression may include ONLY the following operators (listed in the order of their precedence): ; \textasciicircum{} (exponentiation) \(\to\) right associative ; * (multiplication), / (division) \(\to\) left associative ; + (addition), - (subtraction) \(\to\) left associative}
\question{12}{2}{Write a C function bool isDuplicate(char s[]) that detects duplicate parentheses in a given expression using a stack.}
\question{13}{2}{Using the following nested calls, justify why a stack is preferred over a queue in handling function calls. Explain the role of activation records (stack frames).\par void fun1() \{ printf("Inside fun1\textbackslash{}n"); \}\par void fun2() \{ fun1(); printf("Inside fun2\textbackslash{}n"); \}\par int main() \{ fun2(); printf("Inside main\textbackslash{}n"); \}}
\vfill\begin{center}\small End of question paper\end{center}\end{document}